\begin{abstract}
Reasoning models check, revise, decompose, and explore alternatives. Do task-specific patterns in these operations also identify useful interventions? Thought Atlas analyzes 21,803 annotated traces from seven models across eight benchmark conditions. Operation frequencies and temporal placement distinguish conditions beyond trace length, including when temporal controls preserve operation counts. Changing the requested activity on matched problems alters profiles in four of five tested models. Within conditions, the profiles provide no statistically supported outcome-prediction improvement over the specified controls, although coverage is limited for math. On Qwen3.5-9B, interventions constructed by amplifying domain-associated activation contrasts vary markedly in how faithfully they change the measured behavior. Large short-budget gains accompany changes in answer availability and do not survive matched-total-budget elicitation on the tested transfers. Increasing the dose of two contrasts sharply increases literal repetition. On idea generation, math-associated steering increases answer availability while lowering judged originality among paired answers. These results establish reproducible task information in the measured profiles. For the evaluated steering construction, uneven behavioral fidelity and answer production limit what accuracy changes alone reveal about the benefit of the intended behavior.
\end{abstract}
